% CANDIDATE computational note. Not a theorem. Does not edit main.tex.
% XeLaTeX. Continues PC_RAM1. No number below was invented.
\documentclass[11pt]{article}
\usepackage{fontspec}
\usepackage{amsmath,amssymb}
\usepackage[margin=1in]{geometry}
\usepackage{booktabs}
\usepackage{hyperref}
\setmainfont{Latin Modern Roman}
\hypersetup{colorlinks=true,urlcolor=blue,linkcolor=blue}

\title{COMPUTATIONAL: a 7-bit node index, and monomial space of \(F_\star\)\\
\large PC\_RAM2, not a theorem}
\author{Grok Bot, identity \texttt{j15marti}\\
j15marti/pnp, branch \texttt{main}}
\date{4 October 2026, about 12:06 PM PT}

\begin{document}
\maketitle

P versus NP remains open.

This note continues the two negative results in
\texttt{research/agents/grok/2026-10-04\_1157\_GROK\_PC\_RAM1.tex}
(commit \texttt{9076b69827e025e016e74d8639c767d77e9c2fc5},
blob \texttt{122c5553e35730ca8283089700f952134e8595c8}):
the direct RAM charge of \(20\) bits is not \(\log_2 66\) and not \(\log_2 32\),
and characteristic-2 PC degree met resolution width at \(3\) on one frozen
formula \(F_\star\) without touching product-span.
It does not edit root \texttt{main.tex}.
It does not add a workflow.
It does not claim that P versus NP is solved.
Every COMPUTATIONAL figure below was produced by a process run in this pass.

\section*{Standing project facts, not revived}
\begin{itemize}
\item P versus NP remains open.
\item Product-span rank and the compiler into a charged \((T,S,R)\) model remain \textbf{OPEN}.
\item The mask-and-count inequality \(20\neq\log_2 66\) is about one register encoding. It is not, by itself, a theorem that no \(\lceil\log_2 W\rceil\)-bit charge exists.
\end{itemize}

\section*{Baseline, recomputed before any new charge}
Semantics are the sign-round-robin reduced ordered clause-query DAG of
\texttt{experiments/universal-dag/d5\_m8\_exact.cpp}
(blob \texttt{1ca7c97ca45fce506411d246f0f4c499a1352de1}),
specialized to \(n=4\), \(m=8\), as in the ITER1 accounting and in PC\_RAM1.
Input bits are presence of the \(32\) simple signed \(3\)-clauses.
Promise weight is \(8\).
Output \(1\) means the selected clauses falsify every one of the \(16\) assignments.
Terminals are \(0\) (promise-satisfiable), \(1\) (unsatisfiable), and \(2\) (off promise).
A state is a decision node only when the exclude-child and include-child differ.
Nodes at the same level with the same ordered pair of child identifiers are merged.

Recomputed in this pass, and only then used:
\begin{align*}
\text{nonterminal nodes}&=720,\\
\text{peak level width}&=66
\quad\text{(levels \(15\) and \(16\) both have width \(66\))},\\
\log_2 66&=6.044394119358453,\\
\lceil\log_2 66\rceil&=7.
\end{align*}
At level \(16\), live selected-counts are \(k\in\{2,3,4,5,6\}\) with sizes
\(1+16+32+16+1=66\).
Cell \((t,k)=(16,4)\): \(K=32\) decision identities.
All \(\binom{32}{2}=496\) pairs are separated by some common weight-\(4\) continuation of the remaining \(16\) clauses.
So \(\log_2 K=5\).
The baseline matches the PC\_RAM1 witness. The pass did not stop.

On every level, the reachable decision states \((k,\mathrm{cov})\) are in bijection with the reduced nodes at that level:
no decision identity carries two coverages or two selected-counts.
At level \(16\) the \(66\) coverages are pairwise distinct, so \(k\) is a function of the coverage there.

\section*{Named split of the 20-bit charge}
The direct decider stores a coverage mask of the \(2^4=16\) assignments and a selected-count in \(\{0,\ldots,8\}\).
Program counter excluded, as in PC\_RAM1:
\begin{align*}
S_{\mathrm{mask}}&=16,\\
S_k&=\lceil\log_2 9\rceil=4,\\
S_{\mathrm{mutable}}&=20.
\end{align*}
The clause-index program counter, if charged, is \(\lceil\log_2 33\rceil=6\), and \(16+4+6=26\).
Recomputed comparison, unchanged:
\[
20\neq\log_2 66,\qquad 20\neq\log_2 32=5,\qquad 20\neq 7.
\]
What those \(20\) bits actually contain at level \(16\), computed on the \(66\) live coverages:
assignment-bits \(0,1,2,3\) are constantly \(1\);
assignment-bits \(12,13,14,15\) are constantly \(0\);
only assignment-bits \(4\) through \(11\) vary.
That is \(8\) varying mask bits, not \(16\) live ones, and \(2^8=256>66\), so those \(8\) bits are a redundant code for \(66\) states.
The \(4\)-bit selected-count register is not an independent component at this cut: each of the \(66\) coverages determines \(k\).
So \(\log_2 66\) names the live decision classes.
It does not name the constant mask bits, and it does not name the slack in a fixed \(16\)-bit register.
The inequality \(20\neq\log_2 66\) is register width versus the number of live states, not evidence of \(14\) extra semantically distinct bits.

\section*{Alternative charge: 7 mutable bits plus the program}
\textbf{COMPUTATIONAL.}
Number the decision nodes inside each level by increasing creation identifier.
The local index at every level is an integer in \(\{0,\ldots,65\}\) at the peak, hence it fits in \(7=\lceil\log_2 66\rceil\) bits, and \(6\) bits do not suffice because \(66>64\).
The mutable register holds only that index.
The program is the reduced DAG: \(720\) nodes, \(1440\) child edges, of which \(604\) are terminals in \(\{0,1,2\}\).
Every nonterminal child sits at the next clause index (skip count \(0\), maximum jump \(1\)).
The program counter is the current clause index and is not part of the \(7\) bits.
The machine starts at level \(0\), local index \(0\).
On input bit \(b\) it replaces the index by the local index of the \(b\)-child, or halts on a terminal.

If the program counter is charged as well, the sum is \(7+6=13\), which is still not \(20\).

Universe checked: every weight-\(8\) input, population \(\binom{32}{8}=10518300\).
The \(7\)-bit machine matched the mask-union bit on all \(10518300\) inputs
(\(0\) mismatches, \(0\) outputs of the off-promise terminal \(2\)).
Both sides report \(272\) unsatisfiable inputs.
The largest local index that occurred on this universe is \(65\).

The machine therefore answers correctly on the universe that was checked.
The program that makes the \(7\) bits enough is this \(720\)-node table, not a uniform mask-update loop.
No claim is made that a uniform random-access machine with only the \(32\) clause masks as its program can decide the promise problem in \(7\) bits of work memory.
The compiler from an arbitrary SAT algorithm into the charged model remains \textbf{OPEN}.

\section*{\(F_\star\) monomial space, one configuration model}
\(F_\star\) is the weight-\(8\) selection with sign-round-robin indices
\((0,4,8,12,16,20,24,28)\):
all eight sign patterns on \(\{x_2,x_3,x_4\}\), omitting \(x_1\).
Mask union recomputed as \(\mathtt{0xffff}\).
Resolution width recomputed in this pass: minimum width \(3\), width-\(3\) closure size \(27\).

Monomial-space model actually searched.
A configuration is a set of PCR monomials in the twin literals of \(x_1,x_2,x_3,x_4\)
(no variable carries both polarities).
Space is the number of monomials in the configuration.
One step either downloads one of the eight axiom monomials
(each axiom is one degree-\(3\) monomial, the product of the negated clause literals),
erases one monomial,
multiplies one monomial by one literal,
or, when two monomials differ by exactly one variable's polarity, adds their cancellation.
The added monomial is counted before any later erasure.

\textbf{COMPUTATIONAL.}
At space cap \(4\), with all four variables and multiplication, the exhaustive search visited \(775679\) configurations and did not derive the empty monomial.
At space cap \(5\), a refutation that uses only \(x_2,x_3,x_4\), and does not multiply, derives the empty monomial in \(25\) steps:
\(8\) downloads, \(7\) cancellations, \(10\) erasures, peak space \(5\).
Replaying those \(25\) steps found no illegal step.
The final configuration has size \(5\) and contains the empty monomial.
In this model the minimum monomial space of \(F_\star\) is therefore \(5\).

Comparison.
\(5\) exceeds the recomputed resolution width \(3\).
It does not meet that width.
No product-span matrix for \(F_\star\) was computed in this pass.
A repository search for an \(F_\star\) product-span rank returned no note, and
\texttt{experiments/product\_span/README.md} describes a toy affine encoder, not this CNF.
The monomial-space figure is \textbf{incomparable} to the product-span rank.
The connection remains \textbf{untested}. That arrow stays \textbf{OPEN}.

\section*{Status ledger}
\begin{center}
\begin{tabular}{@{}llp{2.6in}@{}}
\toprule
Item & Status & Note \\
\midrule
P versus NP & OPEN & opening sentence \\
baseline \(720\), width \(66\), \(K=32\) & COMPUTATIONAL & matched; pass continued \\
mask \(16\) + count \(4\) \(=20\) & COMPUTATIONAL & PC excluded \\
\(20\neq\log_2 66\), \(20\neq 5\) & COMPUTATIONAL & unchanged inequality \\
level-\(16\) varying mask bits \(=8\) & COMPUTATIONAL & bits \(4\)--\(11\) \\
\(7\)-bit index machine & COMPUTATIONAL & correct on \(10518300\) inputs \\
\(7+6=13\) if PC is charged & COMPUTATIONAL & still not \(20\) \\
compiler / uniform \(7\)-bit RAM & OPEN & program here is the DAG \\
\(F_\star\) resolution width \(3\), closure \(27\) & COMPUTATIONAL & recomputed \\
\(F_\star\) monomial space \(5\) & COMPUTATIONAL & model stated above \\
monomial space vs width & exceeds & \(5>3\), this formula only \\
monomial space vs product-span & incomparable & no matrix; untested \\
product-span and compiler & OPEN & \\
\bottomrule
\end{tabular}
\end{center}

\section*{What this note does not say}
It does not solve P versus NP.
It does not say that every SAT algorithm has mutable space \(\lceil\log_2 W\rceil\).
It does not replace the \(20\)-bit mask charge by a theorem; that charge is still \(20\) bits and still not \(\log_2 66\).
It does not identify monomial space \(5\) with a product-span rank.
The two load-bearing arrows remain OPEN.

\section*{Plain-language summary}
P versus NP is still open.
The ordered \(D(4,8)\) diagram was rebuilt and still has \(720\) nodes, level width \(66\), and a tight cell of \(32\).
The old \(20\)-bit charge splits into a \(16\)-bit mask and a \(4\)-bit selected-count.
At the widest level only \(8\) mask bits vary, and the selected-count is determined by the mask, so the extra bits are register slack.
A different charge stores only the \(7\)-bit index of the current node, with the \(720\)-node diagram as the program, and it matched the true answer on all \(10518300\) weight-\(8\) inputs.
That is a finite simulation for this one program, not a compiler theorem.
On the same frozen \(8\)-clause formula as before, a small monomial-memory refutation needs \(5\) monomials in the model that was searched, which is more than resolution width \(3\), and there is still no product-span matrix to compare it with.

% END_OF_VERIFIED_REPORT::GROK::2026-10-04_1206_PC_RAM2
\end{document}
